% ============================================================
% tick_scheduler.tex
% Section 4: The Tick Scheduler
% Corresponds to: src/core/TickScheduler.py
% Experiments: exp_05_observer_clock.py, exp_07_clock_conservation.py
% ============================================================

\label{sec:scheduler}

The previous section defined a single causal session: a U(1) phase
oscillator on the $\Tdiamond$ lattice with its own tick counter, evolving
under the bipartite tick rule and obeying the $\unity$ constraint at every
tick.
A real physical situation is rarely one session.
Atoms are at least two; molecules many; macroscopic bodies astronomically
many.
The \emph{tick scheduler} is the architectural component that turns a
collection of independent sessions into the dynamics of an interacting
universe.

The scheduler is conceptually small.
It maintains a list of registered sessions, advances each session by one
tick in some order, and applies an interaction step that exchanges phase
between sessions whose probability supports overlap.
That is the whole machine.
The richness of the framework comes from what follows from those two
operations.
This section sets out the formal apparatus.
The phenomenological consequences --- measurement
(Section~\ref{sec:observer}), gravity
(Section~\ref{sec:gravity}), spontaneous emission
(Section~\ref{subsec:emission}) --- are derived from the scheduler in
later sections.

% ── 4.1 The Macro-Tick ────────────────────────────────────────────────────────
\subsection{The Macro-Tick}
\label{subsec:macro_tick}

A \emph{macro-tick} is one cycle of the scheduler.
For a scheduler with $n$ registered sessions
$\{\mathcal{S}_1, \ldots, \mathcal{S}_n\}$, one macro-tick consists of
three steps:

\begin{enumerate}
    \item \textbf{Order.}  Determine a processing order
    $\sigma \in S_n$ over the sessions according to the active shuffle
    scheme.
    \item \textbf{Tick.}  For each $i$ in the order $\sigma$, advance
    $\mathcal{S}_i$ by one bipartite tick (Section~\ref{sec:phase}) and
    increment its tick counter.  Each session enforces $\unity$
    individually after its own tick.
    \item \textbf{Couple.}  Apply the pairwise phase-exchange step
    (Section~\ref{subsec:binding}) at every node where two sessions have
    significant probability support.  Apply any registered three-session
    couplings (Section~\ref{subsec:emission_coupling}).
\end{enumerate}

The decomposition into \emph{tick} and \emph{couple} is essential.
Each session is independently norm-preserving during the tick step ---
this is what makes the U(1) group structure of
Section~\ref{subsec:phase_oscillator} survive at the multi-session level.
The couple step then redistributes \emph{phase} between sessions at
shared nodes without moving probability across them.
$\unity$ is preserved per session through both steps; what is shared is
the coherent structure of the wavefunction, not the probability mass.

In code, the macro-tick is the \texttt{TickScheduler.advance} method.
The full state of the universe at macro-tick $n$ is the tuple
$(\{\mathcal{S}_i\}, \sigma_n, n)$: the sessions, the ordering used at
this step, and the global cycle count.

% ── 4.2 The Combinatorial Clock Space ─────────────────────────────────────────
\subsection{The Combinatorial Clock Space}
\label{subsec:clock_space}

The scheduler's state space is the set of orderings of its registered
clocks.
For $n$ sessions this is the symmetric group $S_n$, of size
\begin{equation}
    |\Omega_n| = n!.
    \label{eq:state_space_factorial}
\end{equation}
This factorial growth is the algebraic root of the framework's
treatment of time.
It will appear again as the entropy of an observer interaction
(Section~\ref{subsec:irreversibility}), as the load on a region near
mass concentrations (Section~\ref{subsec:density_dilation}), and as the
combinatorial source of the second law
(Section~\ref{subsec:arrow_of_time}).

\paragraph{Shuffle schemes as experimental variables.}
The framework does not fix a single ordering rule.
Four schemes are implemented:

\begin{itemize}
    \item \textsc{Sequential}: process sessions in registration order.
    \item \textsc{Random}: uniform random shuffle each macro-tick.
    \item \textsc{Priority}: lowest tick counter first.
    \item \textsc{Clock-density}: highest local
    $\clockdensity$ first.
\end{itemize}

If the framework is correct as a model of physics, the bulk
phenomenology --- conserved probability, the Dirac equation in the
continuum limit, gravitational attraction --- should be invariant under
the choice of scheme.
Different schemes correspond to different choices of simultaneity for the
discrete substrate, in close analogy to choices of foliation in general
relativity.
A discrete \emph{Lorentz invariance} is the prediction that the
choice does not matter for any observable bulk quantity.

\paragraph{Numerical verification (\texttt{exp\_05}).}
The conserved-probability check is straightforward and was performed in
\texttt{exp\_05}.
Across all three implemented schemes
(\textsc{sequential}, \textsc{random}, \textsc{priority}) and across
session counts $n = 1, 2, 3, \ldots, 6$, the total probability satisfies
$|\sum_i \rho_i - n| < 10^{-15}$ at every macro-tick --- machine precision
conservation independent of ordering.
The deeper question --- whether two observers using different shuffle
schemes can disagree on a measurement outcome --- is treated in
Section~\ref{sec:observer}.
The answer is no: the only quantities that depend on the ordering are
those that depend on \emph{which} cone perturbed which, not on the
final probability distribution.

% ── 4.3 Pairwise Coupling and the Binding Dial ────────────────────────────────
\subsection{Pairwise Coupling and the Binding Dial}
\label{subsec:binding}

Two sessions registered in the same scheduler are not, by themselves,
coupled.
They share a lattice but evolve independently.
Coupling enters through the third step of the macro-tick: at any node
$\mathbf{x}$ where two sessions $\mathcal{S}_i$ and $\mathcal{S}_j$ both
have significant amplitude (above a threshold $\sim 10^{-4}$), their
phases are partially exchanged.

The mechanism is a U(1) rotation in the orbital plane of phase
$\phi_i = \arg(\psi_R^{(i)})$ at each shared node:
\begin{align}
    \phi_i &\to \phi_i + c_{ij}\,(\phi_j - \phi_i), \\
    \phi_j &\to \phi_j + c_{ij}\,(\phi_i - \phi_j),
    \label{eq:phase_mix}
\end{align}
where $c_{ij} \in [0, 1]$ is the binding strength registered for the
pair.
Both spinor components are rotated by the same phase
($\unity$ does not mix $\psi_R$ with $\psi_L$ here), and each session is
re-normalised after the rotation.

The single parameter $c_{ij}$ encodes the entire spectrum of two-session
relationships:

\medskip
\begin{center}
\begin{tabular}{ll}
\hline
$c_{ij}$ & Physical regime \\
\hline
$0$       & Free particles --- no interaction \\
$\sim 0.1$ & Weak --- observer / decoherence regime \\
$\sim 0.5$ & Atomic binding --- joint Arnold tongue \\
$\sim 0.9$ & Strong binding --- composite particle \\
$1$       & Phase lock --- sessions move as a single unit \\
\hline
\end{tabular}
\end{center}
\medskip

The same algebraic step thus produces, in different parameter regimes,
the observer effect of Section~\ref{sec:observer}, the joint hydrogen
resonance of Section~\ref{subsec:exp12}, and the bound-state binding
that defines composite particles.
A neutral composite (a neutron, a hydrogen atom seen from far away) is
two or more sessions with $c \to 1$ \emph{and} opposite RGB/CMY
imbalance: their phase gradients cancel, narrowing the effective causal
cone of the composite below those of its constituents.
This cone narrowing is the lattice mechanism for why neutral matter
behaves more classically than its charged constituents.

% ── 4.4 Three-Session Coupling and Emission ───────────────────────────────────
\subsection{Three-Session Coupling and Emission}
\label{subsec:emission_coupling}

Some interactions cannot be reduced to pairwise phase exchange.
The clearest example is spontaneous emission, which we showed in
Section~\ref{subsec:emission} to be a three-session event: the joint
electron--proton resonance jumps to a lower Arnold tongue, and a third
session (the photon) is created to absorb the displaced amplitude under
$\unity$.
The scheduler accommodates this by maintaining an explicit list of
\emph{emission triplets} $(\mathcal{S}_e, \mathcal{S}_\gamma,
\mathcal{S}_p, r)$, where $r$ is the per-tick drain rate.

The drain step is engineered to respect $\unity$ exactly.
At each macro-tick the electron's tangential phase gradient
$\nabla_\perp\phi_e$ (the radial component is preserved so the Coulomb
attraction continues to act) is rotated by $-r\,k_\text{tang}\cdot s$,
where $s$ is the arc-length coordinate along the orbital plane and
$k_\text{tang}$ is the current orbital wavevector.
The conjugate rotation is applied to the photon, transferring orbital
angular momentum from electron to photon while leaving the per-session
probability density --- and therefore $\unity$ --- untouched.
The \emph{appearance} of dissipation comes from the subsequent
re-normalisation step, which sees a redistributed phase structure and
adjusts the spinor components accordingly.

The three-session coupling is the clearest demonstration that the
scheduler is more than a clock-shuffler.
It is a mechanism that decides, at each tick, which combinations of
sessions are interacting and how.
Pairwise binding handles two-session phenomena; emission triplets
handle three-session events; further $n$-session generalisations
(electroweak vertices, hadronic decays) follow the same pattern of
explicit registration with an $\unity$-preserving update rule.

% ── 4.5 Clock Density and Time Dilation ───────────────────────────────────────
\subsection{Clock Density and Time Dilation}
\label{subsec:density_dilation}

The pairwise step of the macro-tick scales with the number of session
overlaps.
For $n$ sessions with significant pairwise support there are at most
$\binom{n}{2}$ pairs to evaluate; for sessions whose probability supports
all share a common node, every pair contributes.
The scheduler is therefore not a constant-cost machine.
Regions of the lattice where many sessions overlap impose more
bookkeeping per macro-tick than regions where sessions are isolated.

This is the lattice substrate for gravitational time dilation.
A region of high clock density $\clockdensity$ runs more couplings per
macro-tick.
A session embedded in that region executes its own tick rule
unchanged --- but each macro-tick it is processed alongside more
neighbours, with more phase exchanges to settle.
From the perspective of a distant observer, who sees only the global
macro-tick count $n$, the local session has accumulated fewer
``personal'' tick advances per global macro-tick: its clock runs slow.

The continuum-limit derivation of this picture --- with the explicit
identification $\nabla\phi_\text{Newton} \leftrightarrow$ macro-tick
imbalance and the recovery of the Schwarzschild time dilation factor
--- is given in Section~\ref{sec:gravity}.
The point here is architectural: time dilation is not added by hand
through a metric.
It is the unavoidable consequence of running a scheduler whose
per-macro-tick cost is proportional to the local session density.

% ── 4.6 Queue Saturation and the Event Horizon ────────────────────────────────
\subsection{Queue Saturation and the Event Horizon}
\label{subsec:saturation}

The scheduler can register at most one session per Planck volume.
Each session occupies a U(1) oscillator at a node, and the lattice
provides a finite number of nodes per unit volume.
The maximum clock density is therefore
\begin{equation}
    \clockdensity^\text{max} = \ell_P^{-3}
        \approx 2.37 \times 10^{104}\ \text{clocks/m}^3,
    \label{eq:rho_max_scheduler}
\end{equation}
the same bound that appears in Section~\ref{subsec:black_holes}.
At this density the scheduler queue is saturated: every available node
is bound to a session, and the next registration cannot proceed.

Saturation is the lattice definition of an event horizon.
From the outside, sessions in the saturated region cannot accept new
clock registrations and so cannot accept new information; the local
tick rate falls to zero relative to the global macro-tick.
This is not a coordinate singularity but a computational deadlock,
strictly bounded above.
There is no infinite curvature inside the horizon because there is no
infinite density: the substrate runs out of nodes before the divergence
can occur.

\paragraph{Hawking radiation as queue-drop.}
A scheduler that has reached saturation but receives further
registration attempts must do something.
The architecturally simplest response is to occasionally drop the
oldest registered session at the boundary, freeing a node for the
incoming registration.
The drop is stochastic over the boundary and proportional to the
attempted registration rate.
Each drop releases a session into the surrounding lower-density region,
where it propagates outward at $c$ along a geodesic of the clock fluid.
The thermal character of the resulting radiation --- the Hawking
spectrum --- is the thermodynamic fingerprint of a uniform-prior
boundary queue under stochastic drop, and is recovered without
trans-Planckian field modes.

\paragraph{The Hawking temperature from boundary surface gravity.}
The temperature of the queue-drop emission can be fixed to leading
order under one modelling assumption: that the saturation boundary
inherits, as an external metric quantity, the Schwarzschild surface
gravity
\begin{equation}
    \kappa = \frac{c^4}{4 G M},
    \label{eq:surface_gravity}
\end{equation}
where $M$ is the total mass enclosed by the saturated region.
This assumption is the architectural counterpart of identifying the
saturation boundary with a black-hole horizon
(Section~\ref{subsec:black_holes}); the strong-field clock-fluid
derivation that fixes the metric to Schwarzschild form is reserved
for the gravitational-continuum-limit work of
Section~\ref{sec:clock_fluid}.

A boundary with effective surface acceleration $\kappa$ supports an
Unruh-style local thermal spectrum at temperature
$T_\text{local} = \hbar\kappa/(2\pi c\,k_B)$.
Folding the gravitational redshift between the local frame and the
asymptotic observer (which, near the horizon, exactly cancels the
$1/\sqrt{1-r_s/r}$ blueshift of $T_\text{local}$) leaves the
temperature seen at infinity equal to the surface-gravity expression
itself~\cite{hawking1975}:
\begin{equation}
    T_H \;=\; \frac{\hbar\,\kappa}{2\pi\,c\,k_B}
        \;=\; \frac{\hbar\,c^3}{8\pi\,G\,M\,k_B}.
    \label{eq:hawking_T}
\end{equation}
The queue-drop rate at the saturated boundary is therefore set, up
to architectural constants of order unity, by
$\Gamma_\text{drop}\sim \kappa/(2\pi)$: the boundary releases on
average one session per $2\pi/\kappa$ of asymptotic time.

\paragraph{Honest scope of Eq.~\eqref{eq:hawking_T}.}
The lattice contribution is the architectural claim that
\emph{some} thermal emission must occur at the saturated boundary
because the scheduler cannot accept further registrations and the
queue-drop is the architecturally simplest legal response.  The
specific form $T_H = \hbar c^3/(8\pi G M k_B)$ then follows from
Eq.~\eqref{eq:surface_gravity} via the standard Unruh--horizon
relation.  The modelling assumption that fixes the numerical
prefactor is the inheritance of the Schwarzschild $\kappa$ at the
saturation boundary: a strong-field derivation of this
correspondence from the clock-fluid equations is the natural
follow-on calculation.

% ── 4.7 Architectural Irreversibility and the Arrow of Time ───────────────────
\subsection{Architectural Irreversibility and the Arrow of Time}
\label{subsec:arrow_of_time}

The scheduler has a \texttt{register\_session} method.
It does not have a \texttt{remove\_session} method, by design.

This single architectural fact is enough to derive the second law of
thermodynamics for the framework.
A measurement, in the sense of Section~\ref{sec:observer}, is the
registration of a new session that overlaps an existing one.
After the registration the state space has grown from $|\Omega_{n-1}| =
(n-1)!$ to $|\Omega_n| = n!$, an increase by a factor of $n$.
The corresponding entropy of the scheduler state space,
\begin{equation}
    S_\text{scheduler} = k_B \log |\Omega_n| = k_B \log(n!),
    \label{eq:scheduler_entropy}
\end{equation}
strictly increases with each registration, and by Stirling's
approximation grows as $S \sim k_B\,n \log n$ for large $n$.
There is no operation that can decrease $n$.
The arrow of time is not a thermodynamic approximation that becomes
exact in some statistical limit --- it is a structural feature of the
scheduler interface.

The deeper consequence, derived in Section~\ref{subsec:irreversibility},
is that what we ordinarily call the second law is the same statement as
``measurement is irreversible.''
Both follow from the absence of a clock-removal operation.
You cannot un-observe a particle for the same reason you cannot reduce
the entropy of an isolated system: the scheduler does not expose the
operation that would let you.

% ── 4.8 Falsifiable Consequences ──────────────────────────────────────────────
\subsection{Falsifiable Consequences}
\label{subsec:scheduler_predictions}

The scheduler architecture makes one prediction directly that distinguishes
it from any continuum-time formalism.

\paragraph{Minimum time-dilation quantum.}
Each registered session contributes one tick of bookkeeping per
macro-tick.
The minimum increment of gravitational time dilation between two
neighbouring regions is therefore one tick duration:
\begin{equation}
    \delta t_\text{min} = t_\text{tick},
    \label{eq:min_dilation}
\end{equation}
where $t_\text{tick}$ is the lattice tick duration in seconds.
At Planck calibration, $t_\text{tick} = t_P \approx 5.4 \times 10^{-44}$~s
and the prediction is unmeasurable.
At Compton calibration, $t_\text{tick} \approx 8 \times 10^{-21}$~s,
which is within roughly two orders of magnitude of current optical-clock
fractional sensitivity ($\sim 10^{-18}$ over millimetre baselines).
Section~\ref{sec:predictions} develops the experimental geometry that
would convert this predicted floor into a falsifiable null-result test:
if optical clocks separated by a calibrated $\Delta\phi$ in a
controlled potential gradient ever resolve a continuous distribution of
fractional offsets without a discrete cutoff, the framework is wrong.

\paragraph{Discrete corrections in queue dynamics.}
Two further structural predictions follow from queue saturation, but
their numerical form requires the gravitational continuum-limit
derivation of Section~\ref{sec:clock_fluid}.
These are: (i) a Planck-density cutoff to gravitational wave dispersion,
producing frequency-dependent gravitational-wave speeds in regions
approaching $\clockdensity^\text{max}$; and (ii) a queue-drop signature
in Hawking emission spectra at the high-frequency tail, where the
finite-lattice cutoff replaces the trans-Planckian modes of standard
treatments.
Both feed into Section~\ref{sec:predictions}.
